\chapter{Literature Review}

\section{Importance of mint as a medicinal and spice plant}

\subsection{Mint (\textit{Mentha spp.})}
Species of mint are commonly grown as medicinal and aromatic plants because of their economic value and wide use in the food, pharmaceutical, and cosmetic industries \textcite{Zheljazkov2020}. Among them \textit{Mentha piperita}, \textit{M. spicata} and \textit{M. crispata} are of special importance due to their high biomass production and essential oil content.

Mint has been known for its medicinal importance since ancient times. Mint was used to create a repellant stench to ward off snakes in the historical records of Nicander of Colophon \textcite{Silva2020}. This points to the longstanding role of mint in traditional medicine and its importance in modern therapeutic applications.

Mint also has value for its essential oil composition. Peppermint leaves contain about 1.2--3.9\% essential oils, which contain more than 300 identified compounds. The most dominant class is the terpenic class with approximately 52\% monoterpenes and 9\% sesquiterpenes, while other components such as aldehydes (9\%), alcohols (6\%), miscellaneous compounds (8\%), lactones (7\%) and aromatic hydrocarbons (9\%) are presented in smaller proportions \textcite{Silva2020}.

Botanically, mint (\textit{Mentha spp.}) is an aromatic herb belonging to the family Lamiaceae, and is mostly perennial. It is well known for its dual function as a medicinal and culinary plant. Members of the genus \textit{Mentha} have been used since ancient civilizations and still play a significant role in modern medicine, food industries and traditional health systems around the world \textcite{Silva2020}.

\subsection{Medicinal Importance of Mint}
Mint has a long and well-established history of medicinal use, especially for the treatment of digestive disorders. Previous studies indicate that \textit{Mentha} species possess carminative, antispasmodic, antiemetic, antimicrobial, antioxidant, anti-inflammatory, and insecticidal properties. These medicinal properties are largely due to essential oils and phenolic compounds such as menthol, menthone, flavonoids, and phenolic acids.

Peppermint (\textit{Mentha piperita}) is one of the most extensively studied species and is commonly used for gastrointestinal ailments, including indigestion, flatulence, colic, nausea, and irritable bowel syndrome. Mint essential oils have also demonstrated antimicrobial activity against various bacteria and fungi, supporting their traditional use in treating infections and preserving food. Additionally, mint oils and menthol are widely used in pharmaceutical preparations, oral hygiene products, cosmetics, and aromatherapy due to their soothing, cooling, and antiseptic effects.

Historical evidence further highlights the medicinal relevance of mint, showing its use across ancient Babylonian, Greek, Roman, Islamic, and medieval European medical practices. According to \textcite{Silva2020}, ``Mints have been used since ancient Babylon, but it was in Classical Antiquity that their medical uses flourished.''

\subsection{Importance of Mint as a Spice Plant}
In addition to its medicinal value, mint is an important spice plant widely used in culinary practices across different cultures. The leaves, which can be consumed fresh or dried, the primary edible parts of the plant ``Either fresh or dried leaves of mint can be used for culinary purposes. Most often fresh mint is chosen over dried mint.''

Fresh leaves are preferred for their stronger aroma and flavor. Mint leaves are characterized by a warm, sweet, and refreshing taste with a distinctive cooling aftereffect, making them a popular flavoring agent in a wide range of foods and beverages, including sauces, jellies, confectionery, syrups, ice creams, and traditional dishes. Mint plays an important role in Middle Eastern, Indian, European, and African cuisines, where it is valued for both flavor enhancement and digestive benefits. Mint essential oils are also widely used in the food industry as natural flavoring agents and preservatives due to their antimicrobial properties.

\subsection{Economic and Cultural Significance}
The combined medicinal and spice value of mint has made it an economically important 
crop in many countries worldwide. The cultivation and processing of mint, 
both fresh and dried leaves, is vital to agricultural and industrial sectors. 
Its continued use in medicine, food, cosmetics, and perfumery underscores the 
enduring importance of mint as a versatile plant of global significance. 
According to \textcite{Silva2020}, ``They are still presently used for numerous 
purposes, including non-medicinal, which makes them economically relevant herbs.''

\section{Botanical and economic importance of \textit{Mentha piperita}, 
\textit{Mentha spicata}, and \textit{Mentha crispata}}

\begin{itemize}
    \item \textit{Mentha piperita} (peppermint) is one of the most economically important mint species due to its high menthol and menthone content. Menthol provides a cooling sensation and is widely used in pharmaceutical preparations, oral hygiene products, cosmetics, and food flavoring.

    \item \textit{Mentha spicata} (spearmint) contains a high amount of carvone in its essential oil. Because of its mild and sweet flavor, spearmint is extensively used in the perfumery industries, flavoring industry, particularly in chewing gum, toothpaste, beverages, and confectionery products.

    \item \textit{Mentha crispata} (curly mint) is considered a variety of spearmint and is distinguished by its curled or wrinkled leaves. It is commonly used in culinary applications such as salads, herbal teas, and garnishes, and also contributes to essential oil production.
\end{itemize}        

\section{Use of peat moss in indoor and greenhouse mint production}
Peat moss is widely used in greenhouse mint production due to its favorable physical properties; however, its extraction is associated with significant environmental impacts, including carbon emissions and peatland degradation \textcite{Cleary2021}. Peatlands store approximately one-third of the world’s soil carbon, but when they are drained for horticultural use, they shift from carbon sinks to major sources of CO\textsubscript{2} emissions. Degraded peatlands contribute nearly 5\% of global CO\textsubscript{2} emissions, exceeding the combined emissions from international aviation and shipping. 

In addition, peat extraction leads to the destruction of important habitats and threatens biodiversity in sensitive ecosystems such as Canada’s Hudson Bay Lowlands. Sustainability is also a major concern, as peat forms very slowly (approximately 1 mm per year), while extraction removes soil at a much faster rate, making it effectively a non-renewable resource. These environmental concerns have prompted regulatory actions, including the United Kingdom’s plan to phase out peat use in horticulture by 2030 and European Union initiatives aimed at reducing peat dependency.

Consequently, biochar has been proposed as a sustainable alternative growing media component due to its porous structure, chemical stability, and potential to improve substrate properties \textcite{Dumroese2018}. Therefore, despite its horticultural advantages, the environmental impacts of peat moss highlight the urgent need for alternative growing media.

\subsection{Overview and Economic Importance}
According to \textcite{yaghobvand2022evaluation}, species of the genus \textit{Mentha} are of high economic value due to their active and aromatic substances and are used as raw materials in food, cosmetics, health, beverage, and pharmaceutical industries. The increasing demand for high-quality essential oils requires a better understanding of the factors influencing plant growth and physiological traits.

\section{Growth Requirements and Physiology of Mint}

\subsection{Mentha piperita (Peppermint)}

\subsection*{Botanical Characteristics}
\textit{Mentha × piperita} is a rhizomatous, upright perennial herbaceous plant. According to the \textcite{rhs2026mentha}, peppermint is a hybrid, originally treated as a species but now known to be a cross between \textit{Mentha aquatica} (watermint) and \textit{Mentha spicata} (spearmint). It produces upright stems with dark green leaves and reddish stems that have a characteristic warm, spicy scent.

\subsection*{Climate and Environmental Requirements}
Research by \textcite{virzo1980adaptability} demonstrated that height growth of \textit{Mentha piperita} is better at 44\% than 100\% daylight. They further reported that at 14\% daylight, growth is limited, probably by lack of photosynthate, and no flowering occurs. Their transfer experiments between light and shade indicated that peppermint responds very effectively to irradiance also in a late stage of the life cycle with changes of specific leaf area and chlorophyll content.

Regarding temperature, environmental stresses including extreme temperatures can diminish yield and increase susceptibility to diseases. In some climates, as noted by the \textcite{umd2026growing}, summer heat may require partial shade; in subtropical regions, summer heat can reduce plant vigor unless some shade is provided.

According to the \textcite{mobot2026piperita}, the plant requires adequate moisture and adapts well to medium to wet soil conditions.

\subsection*{Physiological Responses to Nutrient Stress}
The relationship between photosynthetic carbon metabolism and essential oil biogenesis in peppermint has also been investigated under manganese stress. Under manganese deficiency, $CO_{2}$ exchange rate, total chlorophyll, total assimilatory area, plant dry weight, and essential oil yield were significantly reduced, whereas chlorophyll a/b ratio, leaf area ratio, and leaf stem ratio significantly increased. In leaves of manganese-deficient plants, $CO_{2}$ incorporation into the primary metabolic pool and essential oil were significantly lower.

Peppermint has likewise been examined under iron deficiency. $CO_{2}$ exchange rate, stomatal conductance, oil content, chlorophyll content, and chlorophyll a/b ratio were significantly affected. Growth parameters such as leaf area ratio, specific leaf weight, and leaf stem ratio were also influenced by iron stress.

\subsection*{Soil and Nutrient Requirements}
According to the \textcite{umd2026growing}, peppermint grows best in rich, moist soils but adapts to a wide range of soil conditions except dry ones. Soil quality, particularly the presence of organic matter and essential nutrients, is crucial for robust growth.

\subsection{Mentha spicata (Spearmint)}

\subsection*{Botanical Characteristics}
According to the \textcite{rhs2026spicata}, spearmint (\textit{Mentha spicata}) produces pointed, slightly crinkled leaves that are lighter green in color than peppermint, with the whole plant having a characteristic sweet smell. It is an aromatic perennial spreading by creeping rhizomes.

\subsection*{Physiological Responses to Water Stress}
A field experiment on spearmint under different irrigation regimes in Southern Italy showed that increasing levels of water stress affected later developing leaves, plant water status, and net photosynthesis from the beginning of stress. Photosynthesis limitation played a critical role during the harvest period. Under severe water stress, the crop restricted water losses by modulating stomatal closure and showed lowered mesophyll conductance. Irrigation treatments did not affect the concentration of organic compounds, while the yield of essential oils was negatively affected by water stress due to reduced crop growth in terms of total and leaf biomass, leaf area index, and crop height.

\subsection*{In Vitro Culture Responses}
Research by \textcite{tisserat2004influence} on physical parameters affecting spearmint cultures in vitro showed that the type of physical support greatly influenced culture growth and morphogenesis. They reported that mint shootlets grown on liquid medium only produced four-fold more fresh weight than shootlets grown on agar medium. Positive correlations occurred between culture vessel capacity and growth, morphogenesis, and carvone levels. An automated plant culture system allowed for the production of approximately fifteen-fold greater biomass compared to culture tubes.

Further research by \textcite{tisserat2008growth} confirmed these findings, demonstrating that carvone was only produced from cultured spearmint shoots and was absent in either roots or callus cultures. They also found that the number of media culture immersions greatly affected growth, morphogenesis, and secondary metabolism, with twelve immersions per day being optimal for growth and morphogenesis.

\subsection*{Comparative Growth Studies}
\textcite{yaghobvand2022evaluation} evaluated five mint species in an aeroponic system under greenhouse conditions. Their results showed that \textit{M. spicata} had the highest leaf-to-stem ratio. Most yield-attributes traits including leaf area, leaf fresh and dry weight, root fresh weight, shoot fresh weight, and total plant dry weight, as well as relative water content and photosynthesis rate, were higher in the second harvest than the first.

\subsection{Mentha crispata (Curly Mint)}

\subsection*{Botanical Characteristics and Taxonomy}
According to the \textcite{rhs2026crispa}, \textit{Mentha spicata} var. \textit{crispa} (curly spearmint) is an aromatic perennial reaching 60 cm in height, spreading by creeping rhizomes. The leaves are bright green, broad, and crinkled. Whorls of tiny, pale purple flowers appear on terminal spikes in summer.

The \textcite{rhs2026crispa} lists several synonyms for this plant, including:
\begin{itemize}
    \item \textit{Mentha spicata} 'Crispa'
    \item \textit{Mentha viridis} var. \textit{crispa}
    \item \textit{Mentha crispa} L. (1753)
\end{itemize}

\subsection*{Cultivation Requirements}
According to the \textcite{rhs2026crispa}, curly mint should be grown in any moist soil including chalk, clay, loam, or sand. Soil should be moist but well-drained, though plants tolerate poorly-drained conditions. The plant requires moist conditions but well-drained soil and should be protected from excess winter wet. It performs best in full sun to partial shade and is suitable for west, south, or east-facing positions.

Regarding growth management, the \textcite{rhs2026crispa} notes that this plant may have the potential to become a nuisance. They recommend restricting root run by planting in deep containers and plunging into soil, or growing in small, contained beds. Plants grown in pots benefit from dividing every few years.

\section{Abiotic stress in mint cultivation with emphasis on salt stress}
Salt stress is a major abiotic constraint affecting plant growth and productivity, particularly in controlled cultivation systems where salinity can accumulate in the root zone, as reported by \textcite{Hasanuzzaman2021}. In medicinal and aromatic plants, salt stress negatively affects growth, photosynthesis, and secondary metabolite production, as described by \textcite{AcostaMotos2017}.

\section{Growth requirements and physiology of mint species}
Mint species require well-aerated substrates, adequate water supply, and balanced nutrient availability to achieve optimal growth and biomass production \textcite{Zheljazkov2020}. Physiological differences among mint species influence their tolerance to environmental stresses \textcite{Singh2022}.

\section{Nutrient solutions in soilless mint production}
Where mint is grown in soilless substrates, nutrition must be supplied largely through the irrigation water, and defined nutrient solutions are used so that all essential elements are present in known concentrations. The most widely used formulation is that of \textcite{Hoagland1950}, published as the revised second edition of Circular 347 of the California Agricultural Experiment Station and developed from earlier work by Hoagland and Snyder in 1933 and by Hoagland and Arnon in 1938.

The solution supplies macronutrients as four salts, namely potassium nitrate, calcium nitrate, ammonium dihydrogen phosphate and magnesium sulfate, together providing approximately 210~ppm nitrogen, 235~ppm potassium, 200~ppm calcium, 31~ppm phosphorus, 64~ppm sulfur and 49~ppm magnesium. Micronutrients are supplied as boric acid, manganese chloride, zinc sulfate, copper sulfate and sodium molybdate. Iron was originally added as ferric tartrate, but is now almost always replaced by a chelated form such as Fe-EDTA, which remains available to plants over a wider pH range.

Because the original formulation was intended for sand and solution culture, in which the nutrient solution is renewed frequently, it is commonly applied at reduced strength in container production with organic substrates, where salts accumulate between irrigations. Many modified versions are in use, differing in the phosphate source, the potassium to calcium ratio and the iron chelate, so that the term ``modified Hoagland solution'' does not denote a single composition.

Hoagland-type solutions are widely used in studies of salinity in mint. \textcite{Morshedloo2025}, for example, supplied \textit{Mentha} $\times$ \textit{gracilis} with Hoagland solution at 25\% strength during early growth, increasing to full strength as the plants matured, before imposing salinity treatments of 0, 50 and 100~mM NaCl.

\section{Salt stress effects on medicinal and spice plants}
Salt stress induces osmotic stress and ion toxicity, leading to reduced growth and yield in medicinal and spice plants. These effects are often more pronounced in container-grown crops due to limited root volume \textcite{Hasanuzzaman2021}.

\section{Physiological responses of plants to salt stress}

\subsection{Chlorophyll content and photosynthesis}
Salinity stress commonly leads to a reduction in chlorophyll content and photosynthetic efficiency as a result of chloroplast damage and stomatal limitation \textcite{Kalaji2021}. These effects directly contribute to reduced biomass accumulation under saline conditions \textcite{MunnsTester2008}.

\subsection{Water relations and osmotic adjustment}
Plants exposed to salt stress experience impaired water uptake and reduced leaf water potential, which can trigger osmotic adjustment mechanisms such as the accumulation of compatible solutes \textcite{AcostaMotos2017}.

\section{Biochar effects on substrate properties and salinity stress}
Biochar application has been shown to improve substrate water-holding capacity, cation exchange capacity, and nutrient retention, which may alleviate the negative effects of salinity stress on plants \textcite{Akhtar2021}. Additionally, biochar can reduce sodium availability in the root zone, thereby mitigating ion toxicity under saline conditions \textcite{Joseph2021}.

\section{Research gap and justification}
The preceding sections establish that mint is an economically important crop, peat moss has significant environmental limitations, biochar offers promise as a sustainable alternative, and salt stress poses a real threat to greenhouse production. However, a critical knowledge gap remains: we do not know whether biochar can help mint plants tolerate salt stress. Specifically, there is limited research on the combined effects of biochar amendment and salt stress on different \textit{Mentha} species. No studies have systematically compared multiple mint species under these conditions to determine whether responses are species-specific or whether biochar's benefits vary across genotypes. 

Understanding these interactions is essential for developing evidence-based recommendations for growers seeking sustainable solutions to salt stress in mint production. This gap provides the rationale for the present study.

\section{Aim and objectives of the study}
The main aim of this study was to evaluate the effects of biochar as a peat alternative on the salt stress responses of three mint species. The specific objectives were to:

\begin{itemize}
    \item Assess growth responses of different mint species under salt stress
    \item Evaluate the influence of biochar-amended growing media on salt stress tolerance
    \item Compare physiological responses among mint species,
    \item Identify interactions between growing media composition and salinity stress.
\end{itemize}